% GENERATED FILE. Edit canonical lessons, then run npm run generate:books.
% canonical-source-hash: fnv1a64:ef77b0990c49b598
% canonical-lessons: GU-C440-jharnu, GU-C440-pravah, GU-C440-sangam, GU-C440-naher, GU-C440-vahavvu

\chapter{Streams and Canals}
\label{ch:gu-streams-and-canals}
\hlchaptermodality{440}

\begin{chapteropening}
\textbf{By the end of this chapter:} \emph{I can talk about running water.}

Complete the last lesson of chapter 440: I can talk about running water.
\end{chapteropening}

\section[\texorpdfstring{\gu{ઝરણું} (jharṇũ)}{ઝરણું (jharṇũ)}]{\gu{ઝરણું} (jharṇũ) — a mountain stream; a spring}

\label{lesson:GU-C440-jharnu}



\begin{quote}
Before the new one: say the Gujarati for a pebble, then the Gujarati for to echo.
\end{quote}

\subsection*{You'll want to know: \gu{ઝરણું}}
\textbf{\gu{ઝરણું}} — \emph{jharṇũ} — \textquotedblleft{}a mountain stream; a spring\textquotedblright{}.

\begin{quote}
\gu{પર્વત} \gu{પરથી} \gu{ઠંડું} \gu{ઝરણું} \gu{વહે} \gu{છે}.
\end{quote}

\begin{quote}
\emph{parvat parthī ṭhanḍũ jharṇũ vahe che.}
\end{quote}

\begin{quote}
A cold stream flows down the mountain.
\end{quote}

A neuter noun.

\begin{cousinweb}
\textbf{\gu{ઝરણું}} (\emph{jharṇũ}) is built on \gu{ઝરવું}, to ooze, to trickle, which CDIAL lists under \emph{*jharati}, a form it reconstructs (CDIAL 5346): water that trickles out of the rock.
\end{cousinweb}

\begin{grammarlens}[title={What you've built}]
One more word for getting around.
\end{grammarlens}

\subsection*{Your turn}
\begin{itemize}
\raggedright
  \item \emph{Say it:} \emph{jharṇũ}
  \item \emph{Say it:} \emph{jharṇũ}, once more
  \item \emph{Say it:} \emph{parvat parthī ṭhanḍũ jharṇũ vahe che.} — A cold stream flows down the mountain.
  \item \emph{Recall:} say the Gujarati for a pebble, then the Gujarati for to echo, then say \emph{jharṇũ} again
\end{itemize}

\begin{tcolorbox}[breakable,colback=teal!4,colframe=teal!35!black,title={Before you move on}]
What does \gu{ઝરણું} mean? (\textquotedblleft{}A mountain stream; a spring\textquotedblright{}.) Say it once more.
\end{tcolorbox}
\section[\texorpdfstring{\gu{પ્રવાહ} (pravāh)}{પ્રવાહ (pravāh)}]{\gu{પ્રવાહ} (pravāh) — a current; a flow}

\label{lesson:GU-C440-pravah}



\begin{quote}
Before the new one: say the Gujarati for to echo, then the Gujarati for a mountain stream.
\end{quote}

\subsection*{You'll want to know: \gu{પ્રવાહ}}
\textbf{\gu{પ્રવાહ}} — \emph{pravāh} — \textquotedblleft{}a current; a flow\textquotedblright{}.

\begin{quote}
\gu{ચોમાસામાં} \gu{નદીનો} \gu{પ્રવાહ} \gu{ઘણો} \gu{ઝડપી} \gu{છે}.
\end{quote}

\begin{quote}
\emph{comāsāmā̃ nadīno pravāh ghaṇo jhaḍpī che.}
\end{quote}

\begin{quote}
In the monsoon the river's current is very fast.
\end{quote}

A masculine noun.

\begin{cousinweb}
\textbf{\gu{પ્રવાહ}} (\emph{pravāh}) is a tatsama, Sanskrit \emph{pravāha}, \textquotedblleft{}a stream, a current\textquotedblright{} (Monier-Williams), built on \emph{vah}, to carry: what carries things along.
\end{cousinweb}

\begin{grammarlens}[title={What you've built}]
One more word for getting around.
\end{grammarlens}

\subsection*{Your turn}
\begin{itemize}
\raggedright
  \item \emph{Say it:} \emph{pravāh}
  \item \emph{Say it:} \emph{pravāh}, once more
  \item \emph{Say it:} \emph{comāsāmā̃ nadīno pravāh ghaṇo jhaḍpī che.} — In the monsoon the river's current is very fast.
  \item \emph{Recall:} say the Gujarati for to echo, then the Gujarati for a mountain stream, then say \emph{pravāh} again
\end{itemize}

\begin{tcolorbox}[breakable,colback=teal!4,colframe=teal!35!black,title={Before you move on}]
What does \gu{પ્રવાહ} mean? (\textquotedblleft{}A current; a flow\textquotedblright{}.) Say it once more.
\end{tcolorbox}
\section[\texorpdfstring{\gu{સંગમ} (sangam)}{સંગમ (sangam)}]{\gu{સંગમ} (sangam) — a confluence (of rivers)}

\label{lesson:GU-C440-sangam}



\begin{quote}
Before the new one: say the Gujarati for a mountain stream, then the Gujarati for a current.
\end{quote}

\subsection*{You'll want to know: \gu{સંગમ}}
\textbf{\gu{સંગમ}} — \emph{sangam} — \textquotedblleft{}a confluence (of rivers)\textquotedblright{}.

\begin{quote}
\gu{બે} \gu{નદીઓના} \gu{સંગમ} \gu{પર} \gu{એક} \gu{મંદિર} \gu{છે}.
\end{quote}

\begin{quote}
\emph{be nadīonā sangam par ek mandir che.}
\end{quote}

\begin{quote}
There is a temple where the two rivers meet.
\end{quote}

A masculine noun.

\begin{cousinweb}
\textbf{\gu{સંગમ}} (\emph{sangam}) is a tatsama, Sanskrit \emph{saṁgama}, \textquotedblleft{}a coming together, a meeting\textquotedblright{} (Monier-Williams); a place where rivers meet is often a holy place.
\end{cousinweb}

\begin{grammarlens}[title={What you've built}]
One more word for getting around.
\end{grammarlens}

\subsection*{Your turn}
\begin{itemize}
\raggedright
  \item \emph{Say it:} \emph{sangam}
  \item \emph{Say it:} \emph{sangam}, once more
  \item \emph{Say it:} \emph{be nadīonā sangam par ek mandir che.} — There is a temple where the two rivers meet.
  \item \emph{Recall:} say the Gujarati for a mountain stream, then the Gujarati for a current, then say \emph{sangam} again
\end{itemize}

\begin{tcolorbox}[breakable,colback=teal!4,colframe=teal!35!black,title={Before you move on}]
What does \gu{સંગમ} mean? (\textquotedblleft{}A confluence (of rivers)\textquotedblright{}.) Say it once more.
\end{tcolorbox}
\section[\texorpdfstring{\gu{નહેર} (naher)}{નહેર (naher)}]{\gu{નહેર} (naher) — a canal}

\label{lesson:GU-C440-naher}



\begin{quote}
Before the new one: say the Gujarati for a current, then the Gujarati for a confluence.
\end{quote}

\subsection*{You'll want to know: \gu{નહેર}}
\textbf{\gu{નહેર}} — \emph{naher} — \textquotedblleft{}a canal\textquotedblright{}.

\begin{quote}
\gu{નહેરનું} \gu{પાણી} \gu{ખેતરો} \gu{સુધી} \gu{પહોંચે} \gu{છે}.
\end{quote}

\begin{quote}
\emph{nahernũ pāṇī khetro sudhī pahõce che.}
\end{quote}

\begin{quote}
The canal water reaches the fields.
\end{quote}

A feminine noun.

\begin{cousinweb}
\textbf{\gu{નહેર}} (\emph{naher}) is usually traced to Arabic \emph{nahr}, \textquotedblleft{}a river, a stream\textquotedblright{}, which came through Persian: a river dug by hand.
\end{cousinweb}

\begin{grammarlens}[title={What you've built}]
One more word for getting around.
\end{grammarlens}

\subsection*{Your turn}
\begin{itemize}
\raggedright
  \item \emph{Say it:} \emph{naher}
  \item \emph{Say it:} \emph{naher}, once more
  \item \emph{Say it:} \emph{nahernũ pāṇī khetro sudhī pahõce che.} — The canal water reaches the fields.
  \item \emph{Recall:} say the Gujarati for a current, then the Gujarati for a confluence, then say \emph{naher} again
\end{itemize}

\begin{tcolorbox}[breakable,colback=teal!4,colframe=teal!35!black,title={Before you move on}]
What does \gu{નહેર} mean? (\textquotedblleft{}A canal\textquotedblright{}.) Say it once more.
\end{tcolorbox}
\section[\texorpdfstring{\gu{વહાવવું} (vahāvvũ)}{વહાવવું (vahāvvũ)}]{\gu{વહાવવું} (vahāvvũ) — to let flow (water); to channel}

\label{lesson:GU-C440-vahavvu}



\begin{quote}
Before the new one: say the Gujarati for a confluence, then the Gujarati for a canal.
\end{quote}

\subsection*{You'll want to know: \gu{વહાવવું}}
\textbf{\gu{વહાવવું}} — \emph{vahāvvũ} — \textquotedblleft{}to let flow (water); to channel\textquotedblright{}.

\begin{quote}
\gu{અમે} \gu{નહેરનું} \gu{પાણી} \gu{ખેતરમાં} \gu{વહાવ્યું}.
\end{quote}

\begin{quote}
\emph{ame nahernũ pāṇī khetarmā̃ vahāvyũ.}
\end{quote}

\begin{quote}
We let the canal water flow into the field.
\end{quote}

Past, with \gu{તેણે} (he or she …): \gu{તેણે} \gu{પાણી} \gu{વહાવ્યું}, \emph{teṇe pāṇī vahāvyũ}.

\begin{cousinweb}
\textbf{\gu{વહાવવું}} (\emph{vahāvvũ}) is the causative partner of \gu{વહેવું}, \textquotedblleft{}to flow\textquotedblright{}, a tadbhava of Sanskrit \emph{vahati}, \textquotedblleft{}carries\textquotedblright{} (CDIAL 11453): to make water flow.
\end{cousinweb}

\begin{grammarlens}[title={What you've built}]
One more word for getting around.
\end{grammarlens}

\subsection*{Your turn}
\begin{itemize}
\raggedright
  \item \emph{Say it:} \emph{vahāvvũ}
  \item \emph{Say it:} \emph{vahāvvũ}, once more
  \item \emph{Say it:} \emph{ame nahernũ pāṇī khetarmā̃ vahāvyũ.} — We let the canal water flow into the field.
  \item \emph{Recall:} say the Gujarati for a confluence, then the Gujarati for a canal, then say \emph{vahāvvũ} again
\end{itemize}

\begin{tcolorbox}[breakable,colback=teal!4,colframe=teal!35!black,title={Before you move on}]
What does \gu{વહાવવું} mean? (\textquotedblleft{}To let flow (water); to channel\textquotedblright{}.) Say it once more.
\end{tcolorbox}
